\documentclass[12pt,a4paper]{article}

\usepackage[utf8]{inputenc}
\usepackage[T1]{fontenc}
\usepackage{lmodern}
\usepackage[magyar]{babel}
\usepackage{geometry}
\usepackage{graphicx}
\usepackage{float}
\usepackage{booktabs}
\usepackage{array}
\usepackage{longtable}
\usepackage{xcolor}
\usepackage[hidelinks]{hyperref}
\geometry{margin=2.5cm}

\setlength{\parskip}{0.45em}
\setlength{\parindent}{0pt}
\renewcommand{\arraystretch}{1.25}

\newcommand{\placeholderfigure}[2]{%
\begin{figure}[H]
    \centering
    \fbox{\parbox[c][6.2cm][c]{0.86\textwidth}{\centering
        \textbf{#1}\\[0.35cm]
        #2\\[0.45cm]
        \emph{Ide illeszthető a saját screenshot.}
    }}
    \caption{#1}
\end{figure}
}

\newcommand{\code}[1]{\texttt{#1}}

\begin{document}

\begin{titlepage}
    \centering

    {\scshape\LARGE Miskolci Egyetem \par}
    \vspace{0.5cm}
    {\scshape\Large Gépészmérnöki és Informatikai Kar \par}
    \vspace{0.5cm}
    {\scshape\large Informatikai Intézet \par}

    \vspace{2.2cm}

    {\Huge\bfseries Virtual Gallery -- Interactive Museum Room\par}
    \vspace{0.45cm}
    {\Large\bfseries 3D grafikai beadandó bemutatási jegyzőkönyve\par}

    \vspace{1.6cm}

    {\Large Féléves feladat\par}

    \vfill

    \begin{flushleft}
        \large
        \textbf{Szerző:} Gál Olivér István \\
        \textbf{Neptun-kód:} ISXBF3 \\
        \textbf{Technológia:} C, SDL2, OpenGL, SDL2\_image, OBJ modellbetöltő könyvtár
    \end{flushleft}

    \vspace{1cm}

    \begin{flushright}
        \large
        Miskolc \\
        2026
    \end{flushright}

\end{titlepage}

\tableofcontents
\newpage

\section{Bevezetés}

A beadandó célja egy bejárható, interaktív, esztétikus 3D virtuális galéria elkészítése volt. A program egy múzeumszobát jelenít meg, ahol a felhasználó first-person kamerával mozoghat, külső fájlból betöltött modelleket láthat, textúrázott falak, padló és plafon között járhat, valamint különböző interaktív és vizuális funkciókat próbálhat ki.

A projekt nem egy üres térben forgó modellre épül, hanem egy konkrét koncepcióra: egy virtuális múzeumteremre. Ez azért fontos, mert a grafikai elemek nem különálló technikai demók, hanem egy egységes jelenet részei. A textúrázás, a fények, az animáció, az objektumkijelölés, az árnyékok és az ütközéskezelés mind azt szolgálják, hogy a program használat közben valódi, bejárható tér érzetét adja.

A dokumentáció célja kettős. Egyrészt bemutatja, hogy a program hogyan teljesíti a beadandó minimális követelményeit és milyen többletfunkciókat tartalmaz. Másrészt védési segédletként is használható: a fő fájlok, a program működése és a bemutatás javasolt menete egy helyen szerepelnek, így a védés közben nem kell minden apró kódrészletet külön megnyitni.

\section{A feladat ismertetése}

A feladat egy OpenGL alapú grafikus program elkészítése volt, amelyben a virtuális teret a felhasználó be tudja járni, a program külső modellfájlokat használ, a jelenet textúrázott, fényekkel megvilágított, interaktív elemekkel rendelkezik, és a használatáról futás közben is ad segítséget.

A megoldás témája egy interaktív múzeumszoba. A jelenetben különböző kiállítási tárgyak, festmények, szobrok, talapzatok, lámpák és vitrin jellegű objektumok szerepelnek. A tér mérete és szerkezete szobaszerű: van padló, plafon, falak és meghatározott bejárható terület. A felhasználó WASD billentyűkkel és egérrel mozoghat, kiválaszthat objektumokat, módosíthatja a megvilágítást, kapcsolhatja az animációt, az árnyékokat, valamint a sétáló, emberi szemmagasságú mozgásmódot.

A program fő célja nem csupán az volt, hogy a kötelező funkciók kipipálhatók legyenek, hanem az, hogy ezek egy összefüggő, bemutatható rendszerben jelenjenek meg. Ez védésnél különösen előnyös, mert nem egyenként kell bizonygatni a funkciókat: a felhasználó egyszerűen bejárja a galériát, és közben természetesen látszik a kamerakezelés, a modellbetöltés, a textúrázás, a fénykezelés, az interakció és a renderelési logika.

\section{A megvalósítás célja és fő irányai}

A megvalósítás egyik legfontosabb célja a jól strukturált projektfelépítés volt. A program nem egyetlen hosszú \code{main.c} fájlban található, hanem külön modulokra van bontva. Az alkalmazáslogika, a kamera, a jelenet betöltése, a renderelés, az interakció, a textúrázás, a CSV-kezelés és a súgó külön fájlokba kerültek.

A fejlesztés fő irányai a következők voltak:

\begin{itemize}
    \item bejárható 3D tér létrehozása first-person kamerával,
    \item külső OBJ modellek betöltése külön modellbetöltő könyvtár segítségével,
    \item a jelenet tartalmának adatfájlból, CSV konfiguráció alapján történő kezelése,
    \item textúrázott környezet és textúrázott objektumok megjelenítése,
    \item állítható, több fényforrásra épülő világítás,
    \item időalapú animáció a forgó szobrokhoz,
    \item egérrel történő objektumkijelölés,
    \item stencil buffer alapú outline kiemelés,
    \item átlátszó vitrin, egyszerű árnyékvetítés és ütközéskezelés megvalósítása,
    \item futás közbeni súgó és információs panel biztosítása.
\end{itemize}

A program így egyszerre teljesíti a kötelező grafikai elvárásokat, és több olyan többletfunkciót is tartalmaz, amelyekkel a beadandó védésénél érdemes külön foglalkozni.

\section{Követelmények teljesítése}

Az alábbi táblázat összefoglalja, hogy a beadandó minimális követelményei hogyan jelennek meg a programban.

\begin{longtable}{p{0.27\textwidth} p{0.61\textwidth}}
\toprule
\textbf{Követelmény} & \textbf{Megvalósítás a programban} \\
\midrule
Kamerakezelés & A felhasználó WASD billentyűkkel mozoghat, jobb egérgombbal és egérmozgatással körbenézhet. Szabad mozgásban Q/E billentyűkkel függőlegesen is lehet mozogni. \\
\midrule
Térbeli objektumok & A jelenet több 3D objektumot tartalmaz, például szobrokat, festményeket, talapzatokat, lámpákat és vitrint. A modellek OBJ fájlból töltődnek be. \\
\midrule
Külső modellfájlok betöltése & Az objektumok a \code{scene.csv} fájl alapján kerülnek betöltésre, a tényleges OBJ betöltést a statikusan linkelt OBJ modellbetöltő könyvtár végzi. \\
\midrule
Animáció és interaktivitás & A szobrok időalapú forgó animációt kapnak, amely az R billentyűvel kapcsolható. Kattintással objektumok jelölhetők ki. \\
\midrule
Textúrák & A fal, padló, plafon és a modellek textúrái SDL2\_image segítségével töltődnek be, majd OpenGL textúraként jelennek meg. \\
\midrule
Fények állítása & A fényintenzitás a + és - billentyűkkel módosítható. A változás nem egyszerű képernyőfényerő, hanem az OpenGL fényparamétereit módosítja. \\
\midrule
Használati útmutató & Az F1 billentyű egy súgó overlayt jelenít meg, amely tartalmazza a program kezelését. \\
\midrule
Strukturált felépítés & A program külön modulokra van bontva: alkalmazás, kamera, scene core, scene loader, scene render, interaction, texture, csv, help és utils. \\
\bottomrule
\end{longtable}

A minimális elvárások tehát nem csak névlegesen teljesülnek, hanem a program használata közben közvetlenül bemutathatók.

\section{Megvalósított többletfunkciók}

A projekt több olyan plusz funkciót is tartalmaz, amelyek túlmutatnak az alapkövetelményeken. Ezek közül védésnél különösen érdemes kiemelni az objektumkijelölést, a stencil outline-t, az árnyékokat, az átlátszó vitrint és az ütközéskezelést.

\begin{table}[H]
\centering
\begin{tabular}{p{0.30\textwidth} p{0.58\textwidth}}
\toprule
\textbf{Plusz funkció} & \textbf{Rövid leírás} \\
\midrule
Egérrel történő picking & Bal kattintással kiválasztható a jelenet egyik objektuma. A rendszer sugarat számol a kamera és az egérpozíció alapján, majd a legközelebbi metszett objektumot választja ki. \\
Stencil buffer outline & A kijelölt objektum körül sárgás kiemelés jelenik meg. Ez stencil bufferrel történik, nem egyszerű 2D kerettel. \\
Human mode & A B billentyűvel bekapcsolható sétáló mód, amely emberi szemmagasságot és finom járásérzetet ad. \\
Ütközéskezelés & A kamera nem mehet át indokolatlanul a fontosabb objektumokon és a terem határain. \\
Átlátszó vitrin & A \code{case\_glass} típusú objektum külön átlátszó renderelési passban jelenik meg. \\
Planar shadow & Az objektumok alatt egyszerű, síkra vetített árnyék jelenik meg, amely H billentyűvel kapcsolható. \\
Átméretezhető ablak és fullscreen & Az ablak méretezhető, F11 vagy Alt+Enter segítségével teljes képernyőre váltható. \\
Információs panel & A kijelölt objektum neve és rövid leírása megjelenik a képernyő alsó részén. \\
\bottomrule
\end{tabular}
\caption{A programban megvalósított fontosabb többletfunkciók.}
\end{table}

Ezek a funkciók együtt adják a program igazi karakterét. A beadandó így nem csak technikailag működik, hanem bemutatás közben is látványos: van mire kattintani, van mit kapcsolni, és a néző azonnal látja a különbséget.

\section{A program használata}

A program irányítása egyszerű és védés közben gyorsan bemutatható. A legfontosabb vezérlők a következők:

\begin{table}[H]
\centering
\begin{tabular}{p{0.28\textwidth} p{0.58\textwidth}}
\toprule
\textbf{Billentyű / egér} & \textbf{Funkció} \\
\midrule
W / S & Előre és hátra mozgás \\
A / D & Oldalirányú mozgás \\
Jobb egérgomb + egérmozgatás & Körbenézés first-person nézetben \\
Q / E & Fel és le mozgás szabad módban \\
B & Human mode be- vagy kikapcsolása \\
R & Animáció be- vagy kikapcsolása \\
H & Árnyékok be- vagy kikapcsolása \\
Bal kattintás & Objektum kijelölése \\
+ / - & Fényintenzitás növelése vagy csökkentése \\
F1 & Súgó overlay megjelenítése vagy elrejtése \\
F11 / Alt+Enter & Teljes képernyős mód váltása \\
ESC & Kilépés \\
\bottomrule
\end{tabular}
\caption{A program legfontosabb kezelőbillentyűi.}
\end{table}

\placeholderfigure{A virtuális galéria fő nézete}{Javasolt screenshot: a teljes múzeumterem látszódjon, falakkal, padlóval, objektumokkal és világítással.}

\placeholderfigure{F1 súgó overlay}{Javasolt screenshot: a program futása közben megnyitott súgó, ahol a vezérlés látható.}

\placeholderfigure{Objektumkijelölés és outline}{Javasolt screenshot: egy kijelölt szobor vagy vitrin, körülötte a stencil alapú kiemeléssel és alul az információs panellel.}

\section{A projekt felépítése}

A projekt mappaszerkezete a beadandó szempontjából letisztult. Az alkalmazás fő része az \code{app} könyvtárban található. A \code{src} tartalmazza a forrásfájlokat, az \code{include} a publikus fejléceket, a \code{lib} pedig az előre lefordított OBJ modellbetöltő statikus könyvtárat.

\begin{verbatim}
app/
  Makefile
  src/
    main.c
    app.c
    camera.c
    scene_core.c
    scene_loader.c
    scene_render.c
    scene_interaction.c
    texture.c
    csv.c
    help.c
    utils.c
    scene_internal.h
  include/
    app.h
    camera.h
    scene.h
    texture.h
    csv.h
    help.h
    utils.h
    obj/
      draw.h
      info.h
      load.h
      model.h
      transform.h
  lib/
    libobj.a

demos/
  (órai gyakorlati feladatok)
\end{verbatim}

A \code{demos} mappa az órai gyakorlati feladatokat tartalmazza, a beadandó tényleges futtatása és bemutatása viszont az \code{app} mappából történik.

\subsection{Miért jó ez a szerkezet?}

A moduláris felépítés előnye, hogy a program fő részei elkülönülnek. A kamera nem keveredik a rendereléssel, a CSV beolvasás nem keveredik az ütközéskezeléssel, és a jelenet betöltése sem ugyanabban a fájlban történik, mint a stencil outline vagy az árnyékrajzolás.

Ez védésnél is nagy előny: ha a tanár rákérdez egy funkcióra, gyorsan megmutatható a hozzá tartozó fájl. Például a kamerakezelés a \code{camera.c}, az egérrel történő kijelölés a \code{scene\_interaction.c}, az árnyék és a stencil outline pedig a \code{scene\_render.c} fájlban található.

\section{Adatvezérelt jelenetkezelés}

A jelenet tartalma nem keményen beégetett, repetitív C kódrészekből áll, hanem CSV konfiguráció alapján töltődik be. Ez tisztább és skálázhatóbb megoldás: ha új objektumot kell hozzáadni a teremhez, akkor elsősorban az adatfájlt kell módosítani, nem a programlogikát.

A \code{scene.csv} formátuma:

\begin{verbatim}
type,model,texture,px,py,pz,rx,ry,rz,sx,sy,sz
\end{verbatim}

Az egyes mezők jelentése:

\begin{itemize}
    \item \code{type}: az objektum típusa, például \code{statue}, \code{painting}, \code{pedestal}, \code{lamp} vagy \code{case\_glass},
    \item \code{model}: az OBJ modell fájl elérési útja,
    \item \code{texture}: a textúrafájl elérési útja,
    \item \code{px, py, pz}: pozíció a térben,
    \item \code{rx, ry, rz}: forgatás fokban,
    \item \code{sx, sy, sz}: skálázás.
\end{itemize}

A CSV-alapú jelenetleírás azért erős megoldás, mert a program adatai és a program működése elkülönülnek. Magyarul: nem a kódba van belekalapálva minden tárgy helye. Ez szebb, tisztább, és kevésbé fáj karbantartani.

\section{A fő fájlok szerepe}

\subsection{\code{main.c}}

A \code{main.c} a program belépési pontja. Létrehoz egy \code{App} struktúrát, meghívja az inicializálást, majd futtatja a fő ciklust. A fő ciklus három fő lépésből áll: eseménykezelés, állapotfrissítés és renderelés. Ez a klasszikus grafikus alkalmazás felépítés: input, update, render.

A védésen ezt nagyon egyszerűen lehet elmondani: a \code{main.c} nem akar okos lenni, csak összefogja a program életciklusát. Ez jó jel, mert azt mutatja, hogy a logika ki van szervezve külön modulokba.

\subsection{\code{app.c} és \code{app.h}}

Az \code{app.c} kezeli az alkalmazás magas szintű működését. Itt történik az SDL inicializálása, az OpenGL ablak és context létrehozása, a renderállapotok beállítása, az események feldolgozása, a fullscreen váltás, valamint a fő update és render függvények meghívása.

Fontosabb feladatai:

\begin{itemize}
    \item SDL és SDL2\_image inicializálása,
    \item OpenGL context létrehozása,
    \item alap OpenGL állapotok beállítása,
    \item billentyűzet- és egéresemények kezelése,
    \item F1, R, H, B, +, -, F11 és ESC funkciók feldolgozása,
    \item frame-idő számítása az időalapú animációhoz,
    \item kamera és jelenet frissítése,
    \item jelenet és overlayek kirajzolása,
    \item erőforrások felszabadítása.
\end{itemize}

Az \code{App} struktúra tartalmazza az ablakot, az OpenGL contextet, a kamerát, a jelenetet, az időadatokat, valamint az ablak- és viewportméreteket. Ez csökkenti a globális változók szükségességét, mert az alkalmazásállapot egy helyen, áttekinthető struktúrában van.

\subsection{\code{camera.c} és \code{camera.h}}

A \code{camera.c} felel a first-person kamera mozgatásáért és nézeti transzformációjáért. A kamera pozíciót, forgatást, sebességet, valamint a human mode-hoz tartozó járásfázist és bob offsetet tárol.

A kamera működésének lényege, hogy a mozgás időfüggő. Ez azt jelenti, hogy a kamera nem attól mozog gyorsabban vagy lassabban, hogy éppen hány FPS-sel fut a program, hanem az eltelt idő alapján frissül. Ez fontos minőségi szempont, mert a grafikai beadandóknál gyakori hiba, ha az animáció és a mozgás a kirajzolások számától függ.

A \code{set\_view} függvény állítja be az OpenGL model-view mátrixot a kamera alapján. A \code{clamp\_camera\_to\_room} gondoskodik arról, hogy a kamera a terem határain belül maradjon.

\subsection{\code{scene\_core.c}}

A \code{scene\_core.c} a jelenet alapállapotát és időbeli változásait kezeli. Itt történik a jelenet inicializálása, a közös anyagparaméterek beállítása, a padló, fal és plafon textúráinak betöltése, valamint az animáció frissítése.

Különösen fontos a \code{compute\_model\_bounds} függvény, amely a betöltött modellek csúcsai alapján meghatározza az objektum határait. Ezeket a határokat később a picking és az ütközéskezelés használja.

A \code{change\_light} függvény a fényintenzitást korlátozott tartományban módosítja. Ez azért korrekt megoldás, mert nem engedi, hogy az érték értelmetlenül negatív vagy túl nagy legyen.

\subsection{\code{scene\_loader.c}}

A \code{scene\_loader.c} a CSV-ben leírt objektumokból építi fel a jelenetet. Beolvassa a \code{scene.csv} sorait, létrehozza az \code{Entity} példányokat, betölti az OBJ modelleket, a textúrákat, kiszámolja a bounding adatokat, majd beállítja a szükséges objektumparamétereket.

A szobrok esetében a fájl automatikusan kezeli azt is, hogy bizonyos statue típusok forgathatóak legyenek, mások pedig ne. A talapzatra helyezéshez a program megkeresi a legközelebbi talapzatot, és a szobrot arra igazítja. Ez apróságnak tűnik, de védésen jól mutat: nem csak ``oda van dobva'' a modell, hanem van logika a jelenet összeállítása mögött.

\subsection{\code{scene\_render.c}}

A \code{scene\_render.c} a program egyik legfontosabb fájlja, mert itt történik a látvány nagy része. Ide tartozik az objektumok transzformációja, a fények beállítása, az anyagparaméterek átadása, a terem kirajzolása, az árnyékpass, az átlátszó vitrin kezelése és a stencil outline.

A renderelés több logikai lépésben történik:

\begin{enumerate}
    \item anyag- és fényparaméterek beállítása,
    \item a textúrázott terem kirajzolása,
    \item opcionális planar shadow pass,
    \item átlátszatlan objektumok kirajzolása,
    \item átlátszó objektumok külön kirajzolása,
    \item kijelölt objektum stencil bufferbe írása,
    \item kijelölési outline megjelenítése.
\end{enumerate}

A stencil outline lényege, hogy a kijelölt objektum először stencil értéket ír, majd a program egy kissé felnagyított verziót rajzol úgy, hogy az csak ott látszódjon, ahol az eredeti objektum nincs. Ettől alakul ki a körvonalhatás. Ez látványos plusz funkció, és technikailag is jól védhető.

\subsection{\code{scene\_interaction.c}}

A \code{scene\_interaction.c} tartalmazza a felhasználói interakció fontosabb matematikai és logikai részeit. Ide tartozik az ütközéskezelés, a mátrixinverzió, a vektor- és mátrixműveletek, valamint az egérrel történő picking.

Az objektumkijelölésnél a program a képernyőkoordinátából normalizált eszközkoordinátát számol, majd az inverz MVP mátrix segítségével visszavetíti azt világkoordinátába. Ebből egy picking ray keletkezik. A program ezt a sugarat összeveti az objektumok bounding sphere adataival, és a legközelebbi találatot választja ki.

Az ütközéskezelésnél a kamera egy egyszerű hengeres vagy kapszulaszerű szereplőként viselkedik. A program a fontosabb objektumokat akadályként kezeli, így a felhasználó nem tud teljesen szabadon átsétálni rajtuk. Ez nem teljes fizikai motor, de a beadandó céljához józan, működő és bemutatható megoldás.

\subsection{\code{texture.c}}

A \code{texture.c} feladata a képfájlok betöltése és OpenGL textúrává alakítása. A program SDL2\_image segítségével olvassa be a képet, RGBA formátumra konvertálja, majd \code{glTexImage2D} segítségével feltölti a GPU oldalra.

A textúráknál ismétlés és lineáris szűrés van beállítva, ami a falak, padló és modellek megjelenését szebbé teszi. A külön textúrafüggvény azért hasznos, mert így a betöltési logika nem ismétlődik minden fájlban.

\subsection{\code{csv.c}}

A \code{csv.c} egyszerű, célzott CSV beolvasót tartalmaz a scene konfigurációhoz. A fájl első sorát fejlécként kezeli, majd soronként beolvassa az objektumtípust, modellútvonalat, textúraútvonalat, pozíciót, forgatást és skálázást.

Ez a modul kicsi, de fontos: ennek köszönhető, hogy a jelenet leírása adatfájlba került, nem pedig repetitív C kódba.

\subsection{\code{help.c} és \code{utils.c}}

A \code{help.c} az F1 súgó megjelenítését kezeli. A súgó nem csak konzolra írja ki a vezérlést, hanem képernyőn megjelenő overlayként is használható. Ez a beadandó egyik kötelező pontját teljesíti.

A \code{utils.c} kisebb segédfüggvényeket tartalmaz, például fok-radián konverziót, 2D szövegkirajzolást és kitöltött képernyőpanel rajzolását. Ezek a HUD és az információs panel miatt fontosak.

\section{Renderelési megoldások}

A program OpenGL fixed-function pipeline-ra épül. Ez klasszikus, tanulási szempontból jól követhető megközelítés: a fények, anyagok, textúrák és transzformációk közvetlen OpenGL hívásokkal jelennek meg.

\subsection{Textúrázott terem}

A terem padlója, falai és plafonja nem külön modellként szerepelnek, hanem a program quadokból rajzolja ki őket. Ez egyszerű, hatékony és jól kontrollálható megoldás. A textúrák ismétléssel kerülnek a felületekre, így a terem nem egyszínű doboznak hat.

\subsection{Világítás}

A jelenet több fényforrást használ. A program a \code{lamp} típusú objektumok alapján keresi meg a fények pozícióját, majd ezekhez állítja be az OpenGL fényeket. A + és - billentyűk a fény intenzitását módosítják, amely az ambient, diffuse és specular komponensekre is hat.

Ez fontos különbség a gyenge megoldásokhoz képest: nem a teljes képernyő fényereje változik, hanem a jelenet megvilágítási paraméterei.

\placeholderfigure{Fényintenzitás változtatása}{Javasolt két screenshot: egyik alacsonyabb, másik magasabb fényintenzitással. A különbség a tárgyak és a terem megvilágításán látszódjon.}

\subsection{Planar shadow}

A program egyszerű síkra vetített árnyékot valósít meg. Ez nem teljes, modern shadow mapping, hanem klasszikus planar shadow technika: a program egy árnyékmátrix segítségével a padló síkjára vetíti az objektumok geometriáját, majd sötét, áttetsző alakzatként rajzolja ki.

A nagyobb modelleknél a program proxy árnyékot is használhat, hogy ne kelljen túl részletes geometriát árnyékként kirajzolni. Ez praktikus kompromisszum: a látvány megvan, a teljesítmény pedig nem szenved feleslegesen.

\subsection{Átlátszó vitrin}

Az átlátszó objektumok külön renderelési lépésben jelennek meg. A \code{case\_glass} típusú objektum blendeléssel, kikapcsolt textúrázással, részben áttetsző anyagparaméterekkel rajzolódik ki. Az átlátszó objektumok külön kezelése azért fontos, mert a blending érzékeny a rajzolási sorrendre és a depth buffer használatára.

\placeholderfigure{Átlátszó vitrin és árnyékok}{Javasolt screenshot: a vitrin üvege, a benne lévő tárgy és a padlón látható árnyékszerű hatás.}

\section{Interakció és objektumkijelölés}

Az interakció egyik legerősebb része az egérrel történő objektumkijelölés. A felhasználó bal kattintással kiválaszthat egy kiállítási tárgyat. A kiválasztott objektum körül outline jelenik meg, az alsó információs panel pedig kiírja az objektum nevét és rövid leírását.

A picking működése röviden:

\begin{enumerate}
    \item a program megkapja az egér képernyőkoordinátáját,
    \item ezt átalakítja normalizált koordinátává,
    \item lekéri a projection és model-view mátrixokat,
    \item kiszámolja az inverz MVP mátrixot,
    \item világkoordinátában létrehoz egy sugarat,
    \item a sugarat metszi az objektumok bounding sphere adataival,
    \item a legközelebbi találat lesz a kijelölt objektum.
\end{enumerate}

Ez jó kompromisszum pontosság és egyszerűség között. Nem háromszögszintű picking történik, de a felhasználói élményhez bőven elegendő, és a védésen világosan elmagyarázható.

\section{Kamera és ütközéskezelés}

A kamera first-person jellegű. A mozgás a kamera aktuális szögéhez igazodik, tehát előre mozgásnál nem fix világtengely mentén halad a felhasználó, hanem abba az irányba, amerre néz. Ez sokkal természetesebb bejárást eredményez.

A \code{Human mode} egy külön mozgásmód, amely emberi szemmagasságra állítja a kamerát, letiltja a szabad függőleges repülést, és finom fejmozgás-szerű bob effektet ad a mozgáshoz. Ez apró, de látványos extra, mert a program kevésbé technikai demónak, inkább bejárható térnek érződik.

Az ütközéskezelés két szinten működik:

\begin{itemize}
    \item a kamera nem hagyhatja el a terem határait,
    \item a kamera nem mehet át a fontosabb, collidable objektumokon.
\end{itemize}

A megoldás nem próbál teljes fizikai szimuláció lenni, de a beadandóhoz teljesen korrekt: a cél az, hogy a felhasználó ne sétáljon át indokolatlanul a kiállítási tárgyakon.

\section{Fordítás és futtatás}

A projekt a \code{make} paranccsal fordítható az \code{app} könyvtárból. A Windowsos, MinGW alapú kurzuskörnyezethez a Makefile a következő fő beállításokat használja:

\begin{itemize}
    \item fordító: \code{gcc},
    \item figyelmeztetések: \code{-Wall -Wextra -Wpedantic},
    \item include útvonal: \code{-Iinclude},
    \item statikus OBJ könyvtár: \code{-Llib -lobj},
    \item SDL2, SDL2\_image és OpenGL linkelése.
\end{itemize}

Fordítás:

\begin{verbatim}
cd app
make
\end{verbatim}

Futtatás:

\begin{verbatim}
./museum.exe
\end{verbatim}

Takarítás:

\begin{verbatim}
make clean
\end{verbatim}

A programhoz tartozó asset fájlok külön csomagként kezelhetők. A README alapján az \code{assets} mappát az \code{app/} könyvtárba kell helyezni, az elvárt struktúra pedig:

\begin{verbatim}
app/
  assets/
    config/
    models/
    textures/
\end{verbatim}

Ez megfelel annak az elvárásnak, hogy a nagyobb asset állományok ne feltétlenül a repository részeként legyenek feltöltve, hanem külön linkről legyenek elérhetők.

\section{Értékelés és összegzés}

A beadandó egy jól körülhatárolt témára épülő, működő grafikai alkalmazás. A virtuális galéria koncepciója alkalmas arra, hogy a kötelező elemek természetesen jelenjenek meg benne: a kamera bejárja a teret, az objektumok valódi kiállítási tárgyakként szerepelnek, a textúrázás a terem hangulatát adja, a fények és árnyékok a térbeliséget erősítik, az interakció pedig használhatóbbá teszi a programot.

A projekt erősségei:

\begin{itemize}
    \item konkrét, védhető koncepcióra épül,
    \item több kötelező és plusz grafikai funkciót tartalmaz,
    \item a modellek és a jelenet adatai külső fájlokból érkeznek,
    \item a kód több modulra van bontva,
    \item a fő funkciók védés közben látványosan bemutathatók,
    \item a program nem csak technikai minimumot teljesít, hanem felhasználói élményt is ad.
\end{itemize}

Összességében a Virtual Gallery beadandó alkalmas arra, hogy bemutassa a grafikus programozás alapvető eszközeit: transzformációkat, textúrázást, fényeket, modelltöltést, interaktív kamerát, eseménykezelést és egyszerűbb renderelési effekteket. A plusz funkciók miatt a projekt nem minimumszagú, hanem rendesen összerakott. Védésen a legerősebb stratégia az, ha először a működő jelenetet mutatjuk meg, majd csak utána bontjuk ki a kódot modulonként.

\end{document}
